\documentclass{amsart}
% For dealing with references we use the comment environment
\usepackage{verbatim}
\newenvironment{reference}{\comment}{\endcomment}
%\newenvironment{reference}{}{}
\newenvironment{slogan}{\comment}{\endcomment}
\newenvironment{history}{\comment}{\endcomment}

% For commutative diagrams we use Xy-pic
\usepackage[all]{xy}

% We use 2cell for 2-commutative diagrams.
\xyoption{2cell}
\UseAllTwocells

% We use multicol for the list of chapters between chapters
\usepackage{multicol}

% This is generally recommended for better output
\usepackage{lmodern}
\usepackage[T1]{fontenc}

% For cross-file-references

% Package for hypertext links:
\usepackage{hyperref}

% For any local file, say "hello.tex" you want to link to please

% Theorem environments.
%
\theoremstyle{plain}
\newtheorem{theorem}[subsection]{Theorem}
\newtheorem{proposition}[subsection]{Proposition}
\newtheorem{lemma}[subsection]{Lemma}

\theoremstyle{definition}
\newtheorem{definition}[subsection]{Definition}
\newtheorem{example}[subsection]{Example}
\newtheorem{exercise}[subsection]{Exercise}
\newtheorem{situation}[subsection]{Situation}

\theoremstyle{remark}
\newtheorem{remark}[subsection]{Remark}
\newtheorem{remarks}[subsection]{Remarks}

\numberwithin{equation}{subsection}

% Macros
%
\def\lim{\mathop{\mathrm{lim}}\nolimits}
\def\colim{\mathop{\mathrm{colim}}\nolimits}
\def\Spec{\mathop{\mathrm{Spec}}}
\def\Hom{\mathop{\mathrm{Hom}}\nolimits}
\def\Ext{\mathop{\mathrm{Ext}}\nolimits}
\def\SheafHom{\mathop{\mathcal{H}\!\mathit{om}}\nolimits}
\def\SheafExt{\mathop{\mathcal{E}\!\mathit{xt}}\nolimits}
\def\Sch{\mathit{Sch}}
\def\Mor{\mathop{\mathrm{Mor}}\nolimits}
\def\Ob{\mathop{\mathrm{Ob}}\nolimits}
\def\Sh{\mathop{\mathit{Sh}}\nolimits}
\def\NL{\mathop{N\!L}\nolimits}
\def\CH{\mathop{\mathrm{CH}}\nolimits}
\def\proetale{{pro\text{-}\acute{e}tale}}
\def\etale{{\acute{e}tale}}
\def\QCoh{\mathit{QCoh}}
\def\Ker{\mathop{\mathrm{Ker}}}
\def\Im{\mathop{\mathrm{Im}}}
\def\Coker{\mathop{\mathrm{Coker}}}
\def\Coim{\mathop{\mathrm{Coim}}}

% Boxtimes
%
\DeclareMathSymbol{\boxtimes}{\mathbin}{AMSa}{"02}

%
% Macros for moduli stacks/spaces
%
\def\QCohstack{\mathcal{QC}\!\mathit{oh}}
\def\Cohstack{\mathcal{C}\!\mathit{oh}}
\def\Spacesstack{\mathcal{S}\!\mathit{paces}}
\def\Quotfunctor{\mathrm{Quot}}
\def\Hilbfunctor{\mathrm{Hilb}}
\def\Curvesstack{\mathcal{C}\!\mathit{urves}}
\def\Polarizedstack{\mathcal{P}\!\mathit{olarized}}
\def\Complexesstack{\mathcal{C}\!\mathit{omplexes}}
% \Pic is the operator that assigns to X its picard group, usage \Pic(X)
% \Picardstack_{X/B} denotes the Picard stack of X over B
% \Picardfunctor_{X/B} denotes the Picard functor of X over B
\def\Pic{\mathop{\mathrm{Pic}}\nolimits}
\def\Picardstack{\mathcal{P}\!\mathit{ic}}
\def\Picardfunctor{\mathrm{Pic}}
\def\Deformationcategory{\mathcal{D}\!\mathit{ef}}

\setlength{\emergencystretch}{3em}
\begin{document}
\title[Stacks Project, Tag 0E9Q]{251 --- Fibre-dimension bounds and relative Tor amplitude of a flat relative dualizing complex}
\maketitle
\noindent Copyright (C) 2005--2025 Johan de Jong. Stacks Project authors. Source: \href{https://stacks.math.columbia.edu/tag/0E9Q}{Tag 0E9Q}.

\noindent Editorial difficulty note (not author text). Difficulty H2: Besides cohomological bounds, the proof must establish uniform Tor amplitude against arbitrary modules over the base. It uses arbitrary Noetherian base change followed by square-zero extensions to detect derived tensor with finite modules, then filtered colimits. This passage from geometric fibre bounds to all-module Tor bounds is nonroutine..

\noindent Editorial provenance note (not author text). History: excerpt from revision \texttt{a0444\allowbreak 6e57e\allowbreak c1fbc\allowbreak 252a8\allowbreak 71afc\allowbreak ec775\allowbreak 2fb28\allowbreak 07b14}, dualizing.tex, lines 5081--5129. Reference presentation was adapted; original-author context and core support blocks were retained or added. The mathematical wording is unchanged. Licensed under GNU FDL 1.2 or later; no invariant sections or cover texts. See ../licenses/COPYING. Context blocks reproduce author definitions and supporting results; they are not additional counted problems.

\section{Relative dualizing complexes in the Noetherian case}
\label{section-relative-dualizing-complexes-Noetherian}

\noindent
Let $\varphi : R \to A$ be a finite type homomorphism of
Noetherian rings. Then we define the {\it relative dualizing
complex of $A$ over $R$} as the object
$$
\omega_{A/R}^\bullet = \varphi^!(R)
$$
of $D(A)$. Here $\varphi^!$ is as in
Section \href{https://stacks.math.columbia.edu/tag/0BZI}{section relative dualizing complex algebraic (Tag 0BZI)}.
From the material in that section we see that
$\omega_{A/R}^\bullet$ is well defined up to (non-unique) isomorphism.



\begin{lemma}
\label{lemma-relative-dualizing-trivial-vanishing}
Let $R \to A$ be a finite type homomorphism of Noetherian rings.
Let $\mathfrak q \subset A$ be a prime ideal lying over
$\mathfrak p \subset R$. Then
$$
H^i(\omega_{A/R}^\bullet)_\mathfrak q \not = 0
\Rightarrow - d \leq i
$$
where $d$ is the dimension of the fibre of $\Spec(A) \to \Spec(R)$
over $\mathfrak p$ at the point $\mathfrak q$.
\end{lemma}

\begin{proof}
Choose a factorization $R \to P \to A$ with $P = R[x_1, \ldots, x_n]$
as in Section \href{https://stacks.math.columbia.edu/tag/0BZI}{section relative dualizing complex algebraic (Tag 0BZI)}
so that $\omega_{A/R}^\bullet = R\Hom(A, P)[n]$.
We have to show that $R\Hom(A, P)_\mathfrak q$
has vanishing cohomology in degrees $< n - d$.
By Lemma \href{https://stacks.math.columbia.edu/tag/0A71}{lemma RHom ext (Tag 0A71)} this means we have to
show that $\Ext_P^i(P/I, P)_{\mathfrak r} = 0$ for $i < n - d$
where $\mathfrak r \subset P$ is the prime corresponding to $\mathfrak q$
and $I$ is the kernel of $P \to A$.
We may rewrite this as
$\Ext_{P_\mathfrak r}^i(P_\mathfrak r/IP_\mathfrak r, P_\mathfrak r)$
by More on Algebra, Lemma
\href{https://stacks.math.columbia.edu/tag/087R}{lemma pseudo coherence and base change ext (Tag 087R)}.
Thus we have to show
$$
\text{depth}_{IP_\mathfrak r}(P_\mathfrak r) \geq n - d
$$
by Lemma \href{https://stacks.math.columbia.edu/tag/0AVZ}{lemma depth (Tag 0AVZ)}.
By Lemma \href{https://stacks.math.columbia.edu/tag/0BUY}{lemma depth flat CM (Tag 0BUY)} we have
$$
\text{depth}_{IP_\mathfrak r}(P_\mathfrak r) \geq
\dim((P \otimes_R \kappa(\mathfrak p))_\mathfrak r) -
\dim((P/I \otimes_R \kappa(\mathfrak p))_\mathfrak r)
$$
The two expressions on the right hand side agree by
Algebra, Lemma \href{https://stacks.math.columbia.edu/tag/00P2}{lemma codimension (Tag 00P2)}.
\end{proof}

\begin{lemma}
\label{lemma-relative-dualizing-flat-vanishing}
Let $R \to A$ be a flat finite type homomorphism of Noetherian rings.
Let $\mathfrak q \subset A$ be a prime ideal lying over
$\mathfrak p \subset R$. Then
$$
H^i(\omega_{A/R}^\bullet)_\mathfrak q \not = 0
\Rightarrow - d \leq i \leq 0
$$
where $d$ is the dimension of the fibre of $\Spec(A) \to \Spec(R)$
over $\mathfrak p$ at the point $\mathfrak q$. If all fibres of
$\Spec(A) \to \Spec(R)$ have dimension $\leq d$, then
$\omega_{A/R}^\bullet$ has tor amplitude in $[-d, 0]$
as a complex of $R$-modules.
\end{lemma}

\begin{proof}
The lower bound has been shown in
Lemma \ref{lemma-relative-dualizing-trivial-vanishing}.
Choose a factorization $R \to P \to A$ with $P = R[x_1, \ldots, x_n]$
as in Section \href{https://stacks.math.columbia.edu/tag/0BZI}{section relative dualizing complex algebraic (Tag 0BZI)}
so that $\omega_{A/R}^\bullet = R\Hom(A, P)[n]$.
The upper bound means that $\Ext^i_P(A, P)$ is zero for $i > n$.
This follows from
More on Algebra, Lemma \href{https://stacks.math.columbia.edu/tag/068X}{lemma perfect over polynomial ring (Tag 068X)}
which shows that $A$ is a perfect $P$-module with
tor amplitude in $[-n, 0]$.

\medskip\noindent
Proof of the final statement. Let $R \to R'$ be a ring homomorphism
of Noetherian rings. Set $A' = A \otimes_R R'$. Then
$$
\omega_{A'/R'}^\bullet =
\omega_{A/R}^\bullet \otimes_A^\mathbf{L} A' =
\omega_{A/R}^\bullet \otimes_R^\mathbf{L} R'
$$
The first isomorphism by Lemma \href{https://stacks.math.columbia.edu/tag/0BZV}{lemma base change relative algebraic (Tag 0BZV)}
and the second, which takes place in $D(R')$, by
More on Algebra, Lemma \href{https://stacks.math.columbia.edu/tag/0661}{lemma base change comparison (Tag 0661)}.
By the first part of the proof
(note that the fibres of $\Spec(A') \to \Spec(R')$ have dimension $\leq d$)
we conclude that $\omega_{A/R}^\bullet \otimes_R^\mathbf{L} R'$
has cohomology only in degrees $[-d, 0]$. Taking $R' = R \oplus M$
to be the square zero thickening of $R$ by a finite $R$-module $M$,
we see that $R\Hom(A, P) \otimes_R^\mathbf{L} M$
has cohomology only in the interval $[-d, 0]$ for any finite $R$-module $M$.
Since any $R$-module is a filtered colimit of finite $R$-modules
and since tensor products commute with colimits we conclude.
\end{proof}
\end{document}
